% ==================== 附录 ====================

\section*{附录}
\addcontentsline{toc}{section}{附录}

\subsection*{附录A\quad 支撑材料文件清单}

支撑材料共41个MATLAB程序文件及9个输入数据文件，按问题分4个目录。运行各目录入口程序即可复现全文结果。

\vspace{0.3em}
\noindent\textbf{Q1/（趋势预测，6个程序，无需外部数据）}
\begin{enumerate}[leftmargin=2em,itemsep=0pt]
  \item \texttt{main\_question1.m}（主程序入口）
  \item \texttt{logistic\_fit.m}（Logistic模型拟合，核心）
  \item \texttt{gompertz\_fit.m}（Gompertz模型，稳健性对照）
  \item \texttt{leave\_one\_out\_cv.m}（留一交叉验证）
  \item \texttt{rolling\_origin\_validation.m}（滚动起点验证）
  \item \texttt{sensitivity\_analysis.m}（上限灵敏度分析）
\end{enumerate}

\vspace{0.2em}
\noindent\textbf{Q2/（因素排序，13个程序 + 4个CSV输入）}
\begin{enumerate}[leftmargin=2em,itemsep=0pt]
  \item \texttt{run\_q2\_all.m}（主程序入口）
  \item \texttt{critic\_weights.m}（CRITIC客观赋权，核心）
  \item \texttt{topsis\_rank.m}（TOPSIS排序，核心）
  \item \texttt{rank\_with\_ties.m}（带并列排名）
  \item \texttt{spearman\_rank.m}（斯皮尔曼相关系数）
  \item \texttt{wilson\_lower\_bound.m}（威尔逊下限）
  \item \texttt{load\_q2\_data.m}（读取输入CSV）
  \item \texttt{sensitivity\_analysis\_q2.m}（灵敏度分析）
  \item \texttt{self\_check\_q2.m}（自检）
  \item \texttt{choose\_q2\_chinese\_font.m}（字体配置）
  \item \texttt{plot\_q2\_mechanism.m / \_ranking.m / \_sensitivity.m}（绘图）
\end{enumerate}

\vspace{0.2em}
\noindent\textbf{Q3/（方案优选，12个程序 + 3个CSV输入）}
\begin{enumerate}[leftmargin=2em,itemsep=0pt]
  \item \texttt{run\_q3\_all.m}（主程序入口）
  \item \texttt{build\_q3\_profiles.m}（构建5类画像）
  \item \texttt{build\_q3\_plans.m}（构建候选方案）
  \item \texttt{q3\_apply\_safety\_constraints.m}（指南范围初筛）
  \item \texttt{q3\_topsis\_rank.m}（组内TOPSIS排序，核心）
  \item \texttt{q3\_rank\_with\_ties.m}（带并列排名）
  \item \texttt{load\_q3\_data.m}（读取输入CSV）
  \item \texttt{q3\_sensitivity\_analysis.m}（灵敏度分析）
  \item \texttt{q3\_model\_checks.m}（自检）
  \item \texttt{plot\_q3\_coverage.m / \_profiles.m / \_ranking.m}（绘图）
\end{enumerate}

\vspace{0.2em}
\noindent\textbf{Q4/（情景模拟，10个程序 + 2个CSV输入）}
\begin{enumerate}[leftmargin=2em,itemsep=0pt]
  \item \texttt{run\_q4\_all.m}（主程序入口）
  \item \texttt{load\_q4\_data.m}（读取人口与情景CSV）
  \item \texttt{build\_q4\_baseline.m}（基准负担模型，核心）
  \item \texttt{run\_q4\_scenarios.m}（低中高情景模拟，核心）
  \item \texttt{q4\_uncertainty\_analysis.m}（蒙特卡洛不确定性分析）
  \item \texttt{q4\_model\_checks.m}（自检）
  \item \texttt{plot\_q4\_burden.m / \_dpsir.m / \_scenarios.m / \_sensitivity.m}（绘图）
\end{enumerate}

\subsection*{附录B\quad 核心模型代码}

以下给出问题一中Logistic模型拟合函数的完整代码。其余程序见电子支撑材料。

\begin{lstlisting}[caption={logistic\_fit.m：两参数Logistic拟合}]
function [r, t0, y_fit] = logistic_fit(T, Y, K)
% 拟合固定物理上限K的两参数Logistic模型
%   y(t) = K / (1 + exp(-r*(t-t0)))
%   T: 年份向量, Y: 超重肥胖率(%), K: 固定上限

model = @(p, t) K ./ (1 + exp(-p(1) .* (t - p(2))));
params0 = [0.05, 2018];
lb = [0.001, 1950];
ub = [0.5, 2100];

options = optimoptions('lsqcurvefit', ...
    'Display', 'off', ...
    'MaxFunctionEvaluations', 10000);

params = lsqcurvefit(model, params0, T, Y, lb, ub, options);
r = params(1);
t0 = params(2);
y_fit = model(params, T);
end
\end{lstlisting}

